\documentclass[12pt]{article}
\usepackage[T1]{fontenc}
\usepackage[utf8]{inputenc}
\usepackage{lmodern}
\usepackage{textcomp}
\usepackage{enumitem}
\usepackage{geometry}
\usepackage{graphicx}
\usepackage{amsmath, amssymb}
\geometry{margin=1in}
\begin{document}
\begin{center}
Chemistry - Undergraduate Level Assessment 13 \
Total Marks: 40 \
Answer all questions.
\end{center}

\section*{Multiple Choice (10 marks)}
\begin{enumerate}[label=\arabic*.]
\item Elements with partially filled 4 f or 5 f orbitals include all of the following EXCEPT
\begin{enumerate}[label=(\alph*)]
    \item Cu
    \item Gd
    \item Eu
    \item Am
\end{enumerate}
\item The strongest base in liquid ammonia is
\begin{enumerate}[label=(\alph*)]
    \item NH 3
    \item NH 2-
    \item NH 4+
    \item N$_{2}$H$_{4}$
\end{enumerate}
\item Reduction of D-xylose with NaBH 4 yields a product that is a
\begin{enumerate}[label=(\alph*)]
    \item racemic mixture
    \item single pure enantiomer
    \item mixture of two diastereomers in equal amounts
    \item meso compound
\end{enumerate}
\item Calculate the Larmor frequency for a proton in a magnetic field of 1 T.
\begin{enumerate}[label=(\alph*)]
    \item 23.56 GHz
    \item 42.58 MHz
    \item 74.34 kHz
    \item 13.93 MHz
\end{enumerate}
\item The fact that the infrared absorption frequency of deuterium chloride (DCl) is shifted from that of hydrogen chloride (HCl) is due to the differences in their
\begin{enumerate}[label=(\alph*)]
    \item electron distribution
    \item dipole moment
    \item force constant
    \item reduced mass
\end{enumerate}
\end{enumerate}

\section*{True / False (10 marks)}
\textit{Answer True or False}
\begin{enumerate}[label=\arabic*.]
\setcounter{enumi}{5}
\item True or False: Both paramagnetism and ferromagnetism require the presence of unpaired electrons in their atomic or molecular structures.
\item True or False: The molecular geometry of thionyl chloride, SOCl 2, is tetrahedral because it has four regions of electron density around the sulfur atom.
\item True or False: Most organic compounds do not contain silicon atoms, making tetramethylsilane an ideal reference for 1 H NMR spectroscopy.
\item True or False: Cr$_{3}$+ cannot be used as a paramagnetic quencher due to its specific electronic configuration.
\item True or False: Raman spectroscopy is a light-scattering technique used to analyze molecular vibrations and chemical compositions.
\end{enumerate}

\section*{Long Form Response (20 marks)}
\textit{Answer each question with a clear and complete explanation.}
\begin{enumerate}[label=\arabic*.]
\setcounter{enumi}{10}
\item Discuss the balancing of the redox reaction represented by the equation MnO 4- + I- + H+ \textless{}-\textgreater{} Mn$_{2}$+ + IO 3- + H$_{2}$O, and analyze the stoichiometric relationship between MnO 4- and I-.
\item Discuss the implications of the Heisenberg uncertainty principle for a quantum mechanical particle in the lowest energy level of a one-dimensional box, focusing on the limitations of measuring position and momentum.
\end{enumerate}

\vfill
\noindent\textit{End of Paper}
\end{document}